%!TEX root = ./main.tex

%!TEX root = ../main.tex

\chapter{Einleitung}
Die Einleitung soll den Ausgangspunkt der Arbeit umreißen, in kurzer Form zur Problemstellung hinführen und
das Interesse der lesenden Person für die Arbeit wecken. Allgemeine Einleitung ins Thema, keine Unternehmens-
oder Produktbeschreibungen, Organigramme u.ä., wenn diese nicht direkt zum Thema führen. Ziele
und Vorgehensweise nicht vermischen.

\section{Problemstellung}

\section{Ziel}

\section{Vorgehensweise}

\section{Aufgabenstellung}
Die Fragestellung der Aufgabe ist zu präzisieren. Insbesondere sind das Umfeld, die vorhandenen Randbedingungen
und Betrachtungsgrenzen darzustellen.


%!TEX root = ../main.tex

\chapter{Grundlagen}
Ausgehend von der Aufgabenstellung ist der derzeitige Stand der Technik für die Lösungsfindung zu beschreiben.
Es sind z.B. die Vor- und Nachteile bisheriger Lösungen bzw. fundamentaler Lösungsprinzipien fundiert
von und ggf. anderen Quellen darzulegen.

%!TEX root = ../main.tex

\chapter{Hauptteil}

Der Text soll knapp und klar sein und die wesentlichen Gedanken der Arbeit beinhalten. Ein gewähltes Verfahren
oder ein bestimmter Lösungsweg muss begründet werden. Es ist nicht notwendig, alle Vorversuche
einzeln zu schildern. Bei Versuchen sind Voraussetzungen und Vernachlässigungen sowie die Anordnung,
Leistungsfähigkeit und Messgenauigkeit der Versuchsanordnung anzugeben.
Die Ergebnisse der Arbeit sind unter Berücksichtigung der Voraussetzungen ausführlich zu diskutieren und
mit den bereits bekannten Anschauungen und Erfahrungen zu vergleichen.
21 / 55
Ziel der Arbeit ist es, eindeutige Folgerungen und Richtlinien für die Praxis zu finden.

%!TEX root = ../main.tex

\chapter{Zusammenfassung}
Aufgabenstellung, Vorgehensweise und wesentliche Ergebnisse werden kurz und präzise dargestellt und kritisch
reflektiert. Die Zusammenfassung ist eigenständig verständlich. Länge ca. 1 bis 1,5 Seiten (Problem,
Ziele, Vorgehensweise, Ergebnisse und Ausblick).

\chapter{Ausblick}


%% Beispiele, können problemlos entfernt werden!
\chapter{Beispiele}
Beispiele, können problemlos entfernt werden!
\section{Literatur}
Beispiel Text \cite[S. 10]{testBuch1}
Beispiel Homepage\cite{urlId} \\

\section{Bilder}
\begin{figure}[H]
	\centering
	\includegraphics[width=0.3\linewidth]{resources/images/logo-dhbw}
	\caption{DHBW Logo}
	\label{fig:logo-dhbw}
\end{figure}

\section{Fußnote und Abkürzung}
Fußnote\footnote{Fußnote}, \ac{DHBW}

\section{Tabelle}
\begin{table}[H]
	\centering
	\caption{Tabellenbeispiel}
	\label{tab:example}
	\begin{tabular}{|l|c|r|}
		\hline
		Spalte 1 & Spalte 2 & Spalte 3 \\
		\hline
		Zeile  &  &  \\
		\hline
		& Zeile &  \\
		\hline
		&  & Zeile \\
		\hline
		\multicolumn{2}{|r|}{Verbunden}	& nicht Verbunden \\
		\hline
	\end{tabular}
\end{table}

\section{Skript}
\lstinline[language=c]|printf("Inline Code");|
\lstinputlisting[language=python,caption={Beispiel python Skript},captionpos=t,label=scr:exanple]{resources/example-script.py}
